\documentclass{mwrep}
\usepackage{polski}
\usepackage[polish]{babel}
\usepackage[utf8]{inputenc}
\usepackage{tgbonum}
\usepackage{graphicx}
\usepackage{tabularx}
\usepackage{hyperref}

\makeatletter
\let\ps@closing\ps@plain
\makeatother

\renewcommand{\labelitemi}{$\bullet$}
\setlength{\parindent}{0cm}
\addto\captionspolish{\renewcommand{\tablename}{Tabela}}

\newcommand*{\titleGP}{\begingroup
\centering

{\large Studencka Pracownia Inżynierii Oprogramowania}\\Instytut Informatyki Uniwersytetu Wrocławskiego\par
\vspace*{16\baselineskip}

{\Large Rafał Hirsz, Karol Wieczorek, Krzysztof Wróbel\par}
\vspace*{\baselineskip}

\rule{\textwidth}{1.6pt}\vspace*{-\baselineskip}\vspace*{2pt}
\rule{\textwidth}{0.4pt}\\[\baselineskip]

{\Huge SYMULATOR LOTU}\\[0.2\baselineskip]

\rule{\textwidth}{0.4pt}\vspace*{-\baselineskip}\vspace{3.2pt}
\rule{\textwidth}{1.6pt}\\[\baselineskip]

\scshape
{\huge Specyfikacja wymagań}\par
\vspace*{2\baselineskip}

\begin{figure}[h]
\centering
\includegraphics[width=5\baselineskip]{flightsim-team-logo.pdf}
\end{figure}
\vfill

{\large Wrocław 2014}\par

\pagebreak

\endgroup}

\begin{document}

\thispagestyle{empty}
\titleGP

\begin{center}
\begin{table}[h]
\begin{center}
\caption{Historia zmian dokumentu}\label{T:Zmiany}
\vspace{3ex}
\begin{tabularx}{1\textwidth}{|l|l|l|X|}
\hline
Data & Wersja & Autor & Opis zmian \\ \hline
2013-12-11 & 1.0 & Krzysztof Wróbel & Utworzenie dokumentu \\
2014-01-06 & 1.1 & Krzysztof Wróbel & Poprawki stylistyczne \\
2014-01-21 & 1.2 & Rafał Hirsz & Poprawki stylistyczne \\
\hline
\end{tabularx}
\end{center}
\end{table}
\end{center}

\pagebreak

\tableofcontents

\chapter{Charakterystyka użytkowników}

Symulator jest grą komputerową przeznaczoną zarówno dla użytkowników z~zerową wiedzą o~lotnictwie, jak i~dla profesjonalnych pilotów. W~zależności od stopnia wtajemniczenia gracza w~zagadnienia awioniki oferuje różne funkcje dostosowane do jego umiejętności.

\section{Nowicjusz}

Scenariusz gry dla osoby nieposiadającej wiedzy z~zakresu lotnictwa przewiduje lekcje podstawowych technik pilotażu lekkich samolotów. Lekcje dotyczą startu i~lądowania, lotu poziomego, zakrętów, podchodzenia do lądowania i~kołowania. Dostępne są także podstawowe misje zręczonościowe oraz lot dowolny.

\section{Entuzjasta lotnictwa}

Zakłada się, że entuzjaści lotnictwa chcą maksymalnie zgłębić swoją wiedzę o~pilotażu. Grupie tej jest oferowany pełny cykl artykułów, samouczków i~lekcji z~następujących dziedzin: ogólna wiedza o~samolocie, prawo lotnicze, zasady ruchu lotniczego, osiągi i~planowanie lotu, meteorologia, nawigacja, procedury operacyjne, zasady lotu, łączność oraz ogólne bezpieczeństwo lotu. Przygotowane są misje o~narastającym poziomie trudności. Przewiduje się, że po ukończeniu kursu entuzjaści będą w~stanie wykonać zadania przeznaczone dla pilotów w~realnym świecie.

\section{Pilot w~realnym świecie}

Program z~założenia ma jak najlepiej symulować rzeczywisty lot. Jego użytkownicy z~licencją pilota mogą cieszyć się symulacją rzeczywistego lotu przy nieporównywalnie niższych kosztach od podróży samolotem lub wynajęcia prawdziwego symulatora. Gracze mogą sprawdzić swoje umiejętności w~każdych warunkach pogodowych i~w~każdym miejscu na Ziemi. Dodatkowo mogą być zaaranżowane nieprzewidziane usterki oraz problemy z~pasażerami na pokładzie. Licencjonowany pilot będzie używał wszystkich elementów gry. W~szczególności jest tutaj mowa o~dokładnym zaplanowaniu trasy lotu i~załadowania samolotu, manewrach na lotnisku, kontroli lotów, nawigowaniu, przestrzeganiu procedur lotniczych, używaniu map, schematów oraz list kontrolnych. Dla tej grupy graczy są przygotowane skomplikowane misje przedstawiające sytuacje, z~którymi spotykają się prawdziwi piloci.

\chapter{Wymagania funkcjonalne}

\section{Konfiguracja gry}

\subsection{Konfigurowanie lotu dowolnego}

Przed rozpoczęciem symulacji lotu dowolnego należy określić:
\begin{itemize}
\item używany samolot,
\item miejsce startu,
\item czas startu,
\item warunki pogodowe
\end{itemize}
oraz opcjonalnie:
\begin{itemize}
\item ładunek samolotu,
\item plan lotu,
\item ewentualne usterki.
\end{itemize}

\subsubsection{Wybór samolotu}
Gracz wybiera samolot na podstawie jego wytwórcy (np. Boeing), rodziny (np. 737) i~modelu (np. 737-800). Może wybrać także jego pomalowanie (jeżeli jest dostępne więcej niż jedno). Podczas wyboru maszyny gracz powinien być poinformowany o~jej podstawowych osiągach. Wybór jest przedstawiany graficznie.

\subsubsection{Wybór miejsca startu}
Użytkownik określa lotnisko, z~którego ma się rozpocząć symulacja. Dostępne są następujące kryteria wyszukiwania:
\begin{itemize}
\item państwo (np. Grecja),
\item miasto (np. Londyn),
\item nazwa lotniska (np. Okęcie),
\item kod lotniska IATA (np. KRK),
\item kod lotniska ICAO (np. VHHH).
\end{itemize}
Dodatkowo na obszarze danego lotniska należy dospecyfikować miejsce startu symulacji, np.:
\begin{itemize}
\item pas startowy 18L,
\item pas startowy 36R,
\item bramka 12,
\item stanowisko odladzania,
\item hangar C.
\end{itemize}
Podane miejsca zależą od konkretnego portu lotniczego.

\subsubsection{Wybór czasu startu}
Symulacja wymaga określenia czasu jej startu (tj. daty i~godziny).

\subsubsection{Wybór warunków pogodowych}
Użytkownik ma do wyboru następujące scenariusze pogodowe:
\begin{itemize}
\item predefiniowany -- wybiera się konkretny zestaw zjawisk atmosferycznych (np. mroźny, bezchmurny dzień lub ulewny deszcz),
\item odpowiedni dla miejsca i~pory symulacji -- pogoda jest określana na podstawie klimatu miejsca startu,
\item aktualizowany -- pogoda jest aktualizowana na bieżąco, na podstawie zewnętrznych źródeł.
\end{itemize}

\subsubsection{Określenie ładunku samolotu}
W~skład określenia ładunku samolotu wchodzą:
\begin{itemize}
\item liczba i~rozmieszczenia pasażerów,
\item liczba i~rozmieszczenia bagaży,
\item ilość i~rozmieszczenia ładunku cargo,
\item ilość i~rozmieszczenia paliwa.
\end{itemize}

\subsubsection{Opracowanie planu lotu}
W~celu opracowania planu lotu gracz musi podać następujące informacje:
\begin{itemize}
\item lotnisko docelowe -- wybierane w~podobny sposób jak lotnisko startu,
\item tryb lotu:
\begin{itemize}
\item lot IFR,
\item lot VFR,
\end{itemize}
\item rodzaj nawigacji:
\begin{itemize}
\item na podstawie VOR i~NDB,
\item korytarzami powietrznymi na małej wysokości,
\item korytarzami powietrznymi na dużej wysokości,
\item trasa bezpośrednia.
\end{itemize}
\end{itemize}
Wygenerowany plan lotu można wydrukować.

\subsubsection{Określenie usterek samolotu}
Gracz może wymusić, by podczas symulacji pojawiły się konkretne awarie poniższych typów:
\begin{itemize}
\item awarie przyrządów (np. wysokościomierza, sztucznego horyzontu),
\item awarie systemów (np. elektrycznego, hydraulicznego, kompresji powietrza),
\item awarie radia,
\item awarie silników.
\end{itemize}
Dodatkowo istnieje możliwość spowodowania losowej usterki.

\subsection{Wybór misji}

Użytkownik ma możliwość wyboru konkretnego zadania do wykonania podczas symulacji, tzw. misji. Nazwy misji są posortowane według stopnia trudności. W~podglądzie jest widoczny krótki opis misji i~szacowany czas jej ukończenia. Po wybraniu misji graczowi są przedstawiane szczegóły, mapy i~schematy potrzebne do jej ukończenia. Podczas wyboru misji nie ma możliwości skonfigurowania lotu swobodnego.

\subsection{Szkoła}
W szkole dostępne są lekcje pilotażu w~formie pisemnej (strony HTML) i~symulacji (podobne do misji).

\subsection{Biblioteka}
Biblioteka jest bazą dokumentów potrzebnych pilotom do prowadzenia samolotów. W~bibliotece znajdują się:
\begin{itemize}
\item instrukcje obsługi samolotów,
\item listy kontrolne,
\item schematy lotnisk,
\item schematy podejść,
\item mapy stref przylotniskowych,
\item mapy przestrzeni powietrznych.
\end{itemize}

Dokumenty dotyczące samolotów powinny być wyszukiwane według modelu samolotu, dotyczące lotnisk -- według nazw miejscowości, lotniska lub kodu lotniska, natomiast mapy przestrzeni powietrznych -- według kraju.

\section{Symulacja lotu}

\subsection{Sterowanie}

W tabeli \ref{T:Sterowanie_klawiatura} zestawiono klawisze służące do sterowania grą i~skutki ich naciśnięcia. W~tabeli \ref{T:Sterowanie_joystickiem} przedstawiono zdarzenia towarzyszące danym ruchom joysticka.\\
Dodatkowo jest możliwe sterowanie symulatorem za pomocą myszy. Klikanie konkretnych elementów kokpitu daje takie same efekty, jak użycie ich w~prawdziwym samolocie.

\begin{center}
\begin{table}[h]
\begin{center}
\caption{Sterowanie klawiaturą}\label{T:Sterowanie_klawiatura}
\vspace{3ex}
\begin{tabularx}{1\textwidth}{|l|X|}
\hline
Klawisz(e) & Skutek \\ \hline
\texttt{Esc} & Zakończenie symulacji lub zamknięcie ekranu szczegółowego widoku \\
\texttt{P} & Pauza \\
\texttt{F1} & Ciąg zerowy \\
\texttt{F2} & Zmniejszenie ciągu \\
\texttt{F3} & Zwiększenie ciągu \\
\texttt{F4} & Ciąg maksymalny \\
\texttt{F5} & Schowanie klap i~skrzeli \\
\texttt{F6} & Stopniowe chowanie klap i~skrzeli \\
\texttt{F7} & Stopniowe wysuwanie klap i~skrzeli \\
\texttt{F8} & Całkowite wysunięcie klap i~skrzeli \\
\texttt{F9} & Widok z~kamery z~tyłu samolotu \\
\texttt{F10} & Widok w~dół \\
\texttt{F11} & Widok od przodu samolotu \\
\texttt{F12} & Widok z~kokpitu \\
\texttt{`} & Otwarcie (zamknięcie) okna kontroli lotów \\
\texttt{1} & Otwarcie szczegółowego widoku sztucznego horyzontu \\
\texttt{2} & Otwarcie szczegółowego widoku przepustnic \\
\texttt{3} & Otwarcie szczegółowego widoku GPS \\
\texttt{4} & Otwarcie szczegółowego widoku radia \\
\texttt{5} & Otwarcie szczegółowego widoku panelu górnego \\
\texttt{/} & Wysunięcie (schowanie) spojlerów \\
\texttt{.} & Zaciśnięcie (zwolnienie) hamulców \\
\texttt{Z} & Włączenie (wyłączenie) autopilota \\
\texttt{L} & Włączenie (wyłączenie) świateł lądowania \\
\texttt{G} & Wysunięcie (schowanie) podwozia \\
\texttt{E} & Uruchomienie silników \\
Strzałki & Poruszanie kamerą \\
\hline
\end{tabularx}
\end{center}
\end{table}
\end{center}

\begin{center}
\begin{table}[h]
\begin{center}
\caption{Sterowanie joystickiem}\label{T:Sterowanie_joystickiem}
\vspace{3ex}
\begin{tabularx}{1\textwidth}{|l|X|}
\hline
Zdarzenie & Skutek \\ \hline
Wychylenie do przodu & Opuszczenie krawędzi spływu statecznika poziomego \\
Wychylenie do tyłu & Opuszczenie krawędzi natarcia statecznika poziomego \\
Wychylenie w~prawo & Wychylenie lotek na prawym i~lewym skrzydle, odpowiednio, do dołu i~do góry \\
Wychylenie w~lewo & Wychylenie lotek na prawym i~lewym skrzydle, odpowiednio, do góry i~do dołu \\
Obrót w~prawo & Wychylenie lotki na stateczniku pionowym w~prawo \\
Obrót w~lewo & Wychylenie lotki na stateczniku pionowym w~lewo \\
\hline
\end{tabularx}
\end{center}
\end{table}
\end{center}

\subsection{Kontrola lotów}

Kontrola lotów jest prowadzona w~osobnym oknie dialogowym w~trybie tekstowym. Językiem komunikatów jest język angielski. Graczowi są proponowane wypowiedzi, które po wyborze są przekazywane kontrolerowi. Kontrola lotów jest zaimplementowana zgodnie z~tekstem \emph{Ustawy z~dnia 3 lipca 2002 r. -- Prawo lotnicze} z~późniejszymi poprawkami.

\chapter{Wymagania niefunkcjonalne}

\section{Realizm}
\subsection{Realizm fizyczny}
Kładzie się nacisk na dokładność symulacji dynamiki newtonowskiej. Obliczony tor ruchu bryły sztywnej, na którą działają siły zewnętrzne powinien w~możliwie najmniejszym stopniu różnić się od teoretycznego modelu newtonowskiego.\\
Działanie poszczególnych części samolotu powinno być odwzorowane według najprecyzyjniejszych znanych modeli.
\subsection{Realizm graficzny}
Należy zadbać o~dokładne wyobrażenia samolotów, lotnisk, terenu, naziemnych pomocy nawigacyjnych, budynków, a~także urządzeń \linebreak wewnątrz kokpitu.

\section{Niezawodność}
Zakłada się, że gracze będą chcieli odtwarzać najdłuższe możliwe przy aktualnej technologii loty, dlatego też program powinien być w~stanie pracować przez okres do 48 godzin.

\section{Interfejsy programu}
Aby korzystać ze wszystkich funkcji symulatora, gracz powinien posiadać sprawne łącze internetowe.

\section{Interfejs sprzętowy}
Minimalne wymagania:
\begin{itemize}
\item procesor: 1,33 GHz,
\item RAM: 1 GB,
\item karta graficzna: 128 MB,
\item miejsce na dysku twardym: 4 GB
\end{itemize}
\vspace{1ex}
Zalecane wymagania:
\begin{itemize}
\item procesor: dwurdzeniowy 2,0 GHz,
\item RAM: 4 GB,
\item karta graficzna: 256 MB,
\item miejsce na dysku twardym: 4 GB
\end{itemize}

\end{document}